\subsection{Fortifications}

Armies of the American Civil War made extensive use of Fortifications, including trench lines, enclosed forts and redoubts, blockhouses, etc. Early in the war it was common only to fortify important points the contained a more or less permanent garrison. In the war's later stages (from 1864 on) both armies habitually dug in at every opportunity.

\begin{enumerate}
  \item The location of Fortifications already on the Map at the start of each scenario is listed in the information for each scenario. Also list in each scenario is the number of other Fortification Counters available for each side - those can be Constructed several times each if the player wishes, the number is only the maximum number of Fortification Counters the side can have on the Map at any one time.

  \item Fortification Counters are Constructed normally, but are not destroyed; instead they are immediately removed from the Map whenever left unoccupied by friendly units, or whenever enemy units enter the hex.

  \item Fortification Counters aid the defense only, they have no effect on units attacking out of hexes containing them.

  \item No new Fortification Counters can be constructed in new hexes during any of the 1862 scenarios, although you can replace any in a hex where one was destroyed.

  \item New Fortification Counters can be Constructed in any hex during the 1864 scenarios.

  \item Hex T-14 (Harper's Ferry) can never contain a Fortification counter.
\end{enumerate}